%% ma: L10-B14-0087
%% lop: 10
%% chuong: 5
%% bai: 14
%% dang: 10.14.12
%% loai: DS
%% muc: [NB, NB, NB, TH]
%% dap_an: "SSĐĐ"
%% nguon: "chuyen-de-cac-so-dac-trung-cua-mau-so-lieu-khong-ghep-nhom-toan-10-chuong-trinh-moi.pdf (GV. Phan Nhật Linh) (P05), Bài 03 Dạng 2 Phần II Câu 4"
%% kien_thuc: "Lớp 10 Bài 13 (số trung bình, trung vị, tứ phân vị, mốt), Bài 14 (khoảng biến thiên, khoảng tứ phân vị, phương sai, độ lệch chuẩn, giá trị bất thường)"
%% trang_thai: chinh-thuc
%% ghi_chu: "Tự duyệt theo yêu cầu giáo viên 22/10/2026."
\begin{ex}
Mẫu số liệu sau cho biết sĩ số của $5$ lớp khối 10 tại một trường THPT:
\begin{center}
\begin{tabular}{|c|c|c|c|c|}\hline
$40$ & $45$ & $46$ & $41$ & $43$\\\hline
\end{tabular}
\end{center}
Xét tính đúng sai của các khẳng định sau:
\choiceTF{Khoảng biến thiên của mẫu số liệu là $R=3$}{Trung vị của mẫu số liệu là $Q_2=46$}{\True Số trung bình của mẫu số liệu là $\overline x=43$}{\True Phương sai của mẫu số liệu là $s^2=5{,}2$}
\loigiai{Sắp xếp: $40$; $41$; $43$; $45$; $46$.

a) Sai: $R=46-40=6$.

b) Sai: $Q_2=43$.

c) Đúng: $\overline x=\dfrac{215}5=43$.

d) Đúng: $s^2=\dfrac{9+4+0+4+9}5=5{,}2$.}
\end{ex}
